% \documentclass[type = bachelor]{whu-thesis}
% \documentclass[type = master]{whu-thesis}
\documentclass[type = master]{whu-thesis}
% \documentclass[type = master,showframe]{whu-thesis}
% \documentclass[type = doctor,showframe]{whu-thesis}
% \documentclass[type = bachelor,showframe]{whu-thesis}
\whusetup{
  info = {
    % title      = {武汉大学本科生毕业论文 \LaTeX{} 模版},
    % title*     = {Thesis Template for Wuhan UniversityThesis Template for Wuhan UniversityThesis Template for Wuhan UniversityThesis Template for Wuhan University},
    department = {数学与统计学院},
    % major      = {数学与应用数学},
    author     = {夏大鱼羊},
    student-id = {2021202012345},
    supervisor = {夏大鱼羊},
    supervisor-academic-title = {教授},
    % master-type = professional,
    master-type = academic,
    doctor-cover-withsubject = false,
    doctor-cover-info-spread = false
    % bachelor-english-cover=false
  },
  style = {
    font = default,
    math-font = default,
    cjk-font = windows,
    cjk-fakefont = true,
  }
}
\usepackage{zhlipsum}
\usepackage{lipsum}
\begin{document}


\begin{abstract}
  \zhlipsum[1]
\end{abstract}

\begin{abstract*}
  \lipsum[1]
\end{abstract*}

\tableofcontents


\begin{notation}
  $\alpha$ & 常量 \\
  $\beta$ & 常量 \\
  $\gamma$ & 变量 \\
\end{notation}



% 正文
\mainmatter

\chapter{测试}

\section{文字测试}

这是武汉大学学位论文模版，欢迎使用。

This is Wuhan University thesis template, welcome to use!

\section{公式}

\subsection{算符、希腊字母}

\[\sum\prod\int\iint\alpha\beta\gamma\xi\zeta\eta\epsilon\varepsilon\theta\vartheta
  \phi\varphi\psi\]


\subsection{%
  \texorpdfstring{\texttt{\char`\\mathcal}}{\textbackslash mathcal}、
  \texorpdfstring{\texttt{\char`\\mathscr}}{\textbackslash mathscr} 和
  \texorpdfstring{\texttt{\char`\\mathbb}}{\textbackslash mathbb} 字母表}

\begin{itemize}
  \item \verb|\mathcal|: $\mathcal{ABCDEFGHIJKLMNOPQRSTUVWXYZ}$
  \item \verb|\mathscr|: $\mathscr{ABCDEFGHIJKLMNOPQRSTUVWXYZ}$
  \item \verb|\mathbb|: $\mathbb{ABCDEFGHIJKLMNOPQRSTUVWXYZ}$
\end{itemize}


\subsection{（不）带编号单行公式}

\begin{equation}
  a^2 + b^2 = c^2.
\end{equation}

\[ a^2 + b^2 = c^2.\]

\subsection{（不）带编号多行公式}

\begin{align}
  \text{sum} & = 1 + 2 + \cdots + n \\
             & = \frac12 n(n+1).
\end{align}

\begin{align*}
  \text{sum} & = 1 + 2 + \cdots + n \\
             & = \frac12 n(n+1).
\end{align*}

\subsection{矩阵}

\[\begin{pmatrix}
  a_{11} & a_{22} & a_{33} \\
  a_{21} & a_{22} & a_{23} \\
  a_{31} & a_{32} & a_{33} \\
\end{pmatrix} \quad
\begin{vmatrix}
  a_{11} & a_{22} & a_{33} \\
  a_{21} & a_{22} & a_{23} \\
  a_{31} & a_{32} & a_{33} \\
\end{vmatrix} \quad
\begin{bmatrix}
  a_{11} & a_{22} & a_{33} \\
  a_{21} & a_{22} & a_{23} \\
  a_{31} & a_{32} & a_{33} \\
\end{bmatrix} \quad
\begin{Bmatrix}
  a_{11} & a_{22} & a_{33} \\
  a_{21} & a_{22} & a_{23} \\
  a_{31} & a_{32} & a_{33} \\
\end{Bmatrix}\]

\section{脚注测试}

测试 \footnote{眼看他起朱楼，眼看他宴宾客，眼看他楼塌了。这青苔碧瓦堆，俺曾睡风流觉，将五十年兴亡看饱。
金粉未消亡，闻得六朝香，满天涯烟草断人肠。怕催花信紧，风风雨雨，误了春光。}

测试 \footnote[3]{君不见，左纳言，右纳史，朝承恩，暮赐死。行路难，不在水，不在山，只在人情反覆间！}


\section{引用测试}

\subsection{参考文献}

测试%\cite{whu-bachelor:1,whu-bachelor:2,whu-bachelor:3,whu-bachelor:5,whu-bachelor:7}

测试%\cite*{whu-bachelor:1,whu-bachelor:2,whu-bachelor:3,whu-bachelor:7}



\section{图表测试}

\begin{figure}[ht]
  \centering
  \includegraphics[width = 5cm]{whu-logo.pdf}
  \caption{武汉大学校徽}
  \label{fig:武汉大学校徽}
\end{figure}

引用图~\ref{fig:武汉大学校徽}

\begin{table}[ht]
  \centering
  \caption{%
    简单的表格和引用 abc 123 %\cite{whu-bachelor:1}
  }
  \label{table:简单的表格}
  \begin{tabular}{cc}
    \hline
    a & b \\ \hline
    c & d \\ \hline
    测试 & 文本 \\ \hline
  \end{tabular}
\end{table}

引用表~\ref{table:简单的表格}


\section{已定义好的一些数学定理环境}


\begin{definition}[测度]
  (参见文献xxx) 这是一段文字 $E = m c^2$  （中文括号）和 (西文括号)
\end{definition}

\begin{theorem}
  这是一段文字 $E = m c^2$
\end{theorem}


\begin{proof}
  这是一段文字 $E = m c^2$
\end{proof}

\begin{proof}[定理xx的证明]
  这是一段文字 $E = m c^2$
\end{proof}

\begin{example}
  这是一段文字 $E = m c^2$
\end{example}

\begin{property}
  这是一段文字 $E = m c^2$
\end{property}

\begin{proposition}
  这是一段文字 $E = m c^2$
\end{proposition}

\begin{corollary}
  这是一段文字 $E = m c^2$
\end{corollary}

\begin{lemma}
  这是一段文字 $E = m c^2$
\end{lemma}

\begin{axiom}
  这是一段文字 $E = m c^2$
\end{axiom}

\begin{counterexample}
  这是一段文字 $E = m c^2$
\end{counterexample}

\begin{conjecture}
  这是一段文字 $E = m c^2$
\end{conjecture}

\begin{question}
  这是一段文字 $E = m c^2$
\end{question}

\begin{claim}
  这是一段文字 $E = m c^2$
\end{claim}

\begin{remark}
  这是一段文字 $E = m c^2$
\end{remark}

\begin{theorem}[Banach-Steinhaus]\label{thm:test}
  设 $E$ 是 Banach 空间, $F$ 是赋范空间, $(u_i)_{i\in I}$ 是一族从 $E$ 到 $F$ 的有界线性算子,
  即 $(u_i)_{i\in I}\subset \mathcal{B}(E,F)$. 若对每一点 $x\in E$, 有
  $\sup_{i\in I} \|u_i(x)\|<\infty$, 则
  \begin{equation}\label{eq:test1}
    \sup_{i\in I} \|u_i\| < \infty.
  \end{equation}
\end{theorem}

我想引用定理~\ref{thm:test} 和公式~\ref{eq:test1}


定理括号测试：

\begin{theorem}
  测试
  \begin{enumerate}
    \item 中文（括号）没输入空格的效果
    \item 中文 （括号） 输入空格的效果
    \item 西文(括号)没输入空格的效果
    \item 西文 (括号) 输入空格的效果
  \end{enumerate}
\end{theorem} 


\begin{proof}
  test
  \[
    \int_{0}^{1} x^2 \d x
  \]
\end{proof}

\begin{proof}
  test
  \[
    \int_{0}^{1} x^2 \d x \qedhere
  \]
\end{proof}


\section{字体测试}
字体测试：\\
{宋体} {\heiti 黑体} {\kaishu 楷书} {\fangsong 仿宋}\\
{\rmfamily 罗马字族} {\sffamily 无衬线字族} {\ttfamily 打字机}\\
{\bfseries 粗体} {\itshape 意大利} {\slshape 倾斜}

伪粗体测试：\\
{\bfseries\songti 伪粗体} {\bfseries\kaishu 伪粗体} {\bfseries\heiti 伪粗体} {\bfseries\fangsong 伪粗体} {\bfseries 伪粗体}

伪斜体测试：\\
{\itshape\songti 伪斜体} {\itshape\kaishu 伪斜体} {\itshape\heiti 伪斜体} {\itshape\fangsong 伪斜体} {\itshape 伪斜体}

叠加测试：\\
{\bfseries\itshape 伪粗斜体} {\bfseries\sffamily 伪粗黑体} {\bfseries\ttfamily 伪粗仿宋} {\itshape\sffamily 伪斜黑体} {\itshape\ttfamily 伪斜仿宋}


% 参考文献
% \nocite{*}
% \bibliography{bachelor-refs}


\appendix

\chapter{测试}

\section{公式测试}
\section{公式测试}

\begin{equation}
  a^2 + b^2 = c^2
\end{equation}

\chapter{测试}

\section{公式测试}
\subsection{编号测试}
\subsection{编号测试}
\subsubsection{编号测试}
\subsubsection{编号测试}

\section{公式测试}

\begin{equation}\label{equation:appendix-test}
  a^2 + b^2 = c^2
\end{equation}

公式~\eqref{equation:appendix-test}

\end{document}